% Discussion (compact): implications for cutoff-bounded reference generation.
\section{Discussion}
\label{sec:discussion}

\paragraph{Input boundaries and inference errors.}
On the original text, withholding later material from page generation avoided stated
disclosures, whereas full-text cutoff instructions did not. This supports restricting
later-text access but does not establish safety (Section~\ref{sec:results-main}):
\sys's implied disclosure arose from
misreading eligible evidence\ifnwdecon{}, as did its stated marriage on the treated text,
which a character predicts in book two\fi{}. Input filtering cannot prevent these inference
errors. These runs motivate access-boundary controls alongside retrieval and graph choices;
they do not isolate their causal effects. Checks must also cover delivered surfaces, since
development leaks occurred in page lists, titles and shared navigation
(Appendix~\ref{app:casebook}).

\paragraph{Benefits and costs of a linked graph.}
For this familiar five-volume work fitting in one context, graph extraction recovered fewer
gold facts than direct prompting, while prefix retrieval achieved higher recall with about
one tenth of the input tokens. The graph provides deterministic structured content, reusable
cutoff filtering and broader linked coverage, whose value to readers remains unmeasured.
The costs include misplaced claims on relationship pages, insufficient citations and
identities shaped by later volumes (Tables~\ref{tab:kinds} and~\ref{tab:identity}).
For character-page generation in this pilot, the measured tradeoffs favor prefix retrieval.

\paragraph{Design implications.}
Validating entity groupings, preserving speaker and modality, and checking final assertions
in page context against their attached citations are repair priorities whose effectiveness
remains untested. They address distinctions that volume filtering alone cannot verify;
simply publishing more content does not establish reader utility.
